\chapter{Manage Catalog as a Service}

The end-to-end path works, but the catalog still depends on a terminal running
\texttt{bin/rails server}. We will put that same Puma-backed Rails process
behind FreeBSD's service interface, then tell Blackship when to start and stop
it. The service remains disabled in \texttt{rc.conf}: Blackship is the one
lifecycle owner.

This exercise changes only catalog. Keeping store in its foreground console
lets us prove the service pattern at one boundary before repeating it.

\section{Exercise 18A: Create the catalog launcher}

In the catalog console, stop the foreground server with \texttt{Ctrl-C}. If the
prompt belongs to the \texttt{catalog} user, return to \texttt{root}:

\begin{lstlisting}[language=sh]
exit
\end{lstlisting}

Create \texttt{/usr/local/sbin/catalog-server}:

\begin{lstlisting}[language=sh]
vi /usr/local/sbin/catalog-server
\end{lstlisting}

Enter:

\begin{lstlisting}[language=sh]
#!/bin/sh
set -eu

cd /usr/local/www/catalog
exec /usr/local/bin/ruby bin/rails server \
  --environment development \
  --binding 10.17.89.24 \
  --port 3000
\end{lstlisting}

Make it executable and validate its shell syntax:

\begin{lstlisting}[language=sh]
chmod 0555 /usr/local/sbin/catalog-server
sh -n /usr/local/sbin/catalog-server
\end{lstlisting}

The launcher deliberately uses the application's \texttt{bin/rails}; that
script loads the application's bundle and starts Puma. Running it through
FreeBSD's \texttt{daemon(8)} in the next step supplies detachment, supervision,
logging, and a stable supervisor PID.

\section{Exercise 18B: Add an rc.d service}

Create \texttt{/usr/local/etc/rc.d/catalog}:

\begin{lstlisting}[language=sh]
vi /usr/local/etc/rc.d/catalog
\end{lstlisting}

Enter:

\begin{lstlisting}[language=sh]
#!/bin/sh

# PROVIDE: catalog
# REQUIRE: NETWORKING
# KEYWORD: shutdown

. /etc/rc.subr

name="catalog"
desc="Rails catalog API"
rcvar="catalog_enable"

load_rc_config "$name"

: ${catalog_enable:="NO"}
: ${catalog_logfile:="/var/log/catalog.log"}

pidfile="/var/run/${name}.pid"
required_files="/usr/local/sbin/catalog-server"
command="/usr/sbin/daemon"
procname="/usr/sbin/daemon"
command_args="-f -r -P ${pidfile} -u catalog \
  -o ${catalog_logfile} \
  /usr/local/sbin/catalog-server"

run_rc_command "$1"
\end{lstlisting}

Validate and install its executable mode:

\begin{lstlisting}[language=sh]
sh -n /usr/local/etc/rc.d/catalog
chmod 0555 /usr/local/etc/rc.d/catalog
\end{lstlisting}

Do not run \texttt{sysrc catalog\_enable=YES}. The manifest hook will use the
explicit \texttt{onestart} and \texttt{onestop} actions, which allow Blackship
to own the process lifecycle without also asking \texttt{/etc/rc} to start it.

Start the new service once and inspect it:

\begin{lstlisting}[language=sh]
service catalog onestart
service catalog onestatus
sockstat -4 -l | grep ':3000'
tail -n 20 /var/log/catalog.log
\end{lstlisting}

Expect the status command to report that catalog is running and
\texttt{sockstat} to show Ruby listening on \texttt{10.17.89.24:3000}.

From the FreeBSD host, confirm that the API still responds:

\begin{lstlisting}[language=sh]
fetch -qo - http://10.17.89.24:3000/products
\end{lstlisting}

\section{Exercise 18C: Declare the lifecycle in Blackship}

Open the \texttt{catalog} jail's \textbf{Manifest} editor in BSDRunner. Add this
line alongside its existing \texttt{name}, \texttt{release}, and \texttt{build}
settings:

\begin{lstlisting}
depends_on = ["database"]
\end{lstlisting}

Append the following tables to the fragment:

\begin{lstlisting}
[jails.healthcheck]
enabled = true

[[jails.healthcheck.checks]]
name = "catalog-api"
command = "fetch -qo /dev/null http://10.17.89.24:3000/products"
target = "jail"
start_period = 30
interval = 30
timeout = 5
retries = 3

[[jails.hooks]]
phase = "post_start"
target = "jail"
command = "/usr/sbin/service"
args = ["catalog", "onestart"]
on_failure = "abort"

[[jails.hooks]]
phase = "pre_stop"
target = "jail"
command = "/usr/sbin/service"
args = ["catalog", "onestop"]
on_failure = "continue"
\end{lstlisting}

Run \textbf{Dry run} first. Review the dependency, health-check, and hook
changes, then choose \textbf{Apply}. Applying the manifest records the desired
lifecycle but does not replay a \texttt{post\_start} hook against an already
running jail.

Use BSDRunner to restart \texttt{catalog}. Blackship should run the pre-stop
hook, restart the jail, and run the post-start hook. Wait for the catalog health
check to become healthy, then repeat the host request:

\begin{lstlisting}[language=sh]
fetch -qo - http://10.17.89.24:3000/products
\end{lstlisting}

Finally, reload the externally published store page. It should still show
\texttt{FreeBSD Handbook} and \texttt{\$49.99} even though no terminal owns the
catalog process now.

\section{Checkpoint}

Stop after catalog becomes healthy following the Blackship restart and both
requests succeed. The catalog console may now be closed. The next exercise will
apply the same service pattern to store, then the whole application can survive
managed jail restarts without foreground Rails consoles.
